% Introduction
\section{Introduction}

Is a ``truthful'' LLM truthful in all the same ways? A model achieving 95\% on MMLU can fail 40\% of TruthfulQA questions---and our analysis of 50 diverse language models reveals this is the norm, not the exception. Four models in our study demonstrate the disconnect starkly: high scores on general knowledge retrieval (MMLU) yet poor performance on misconception resistance (TruthfulQA), suggesting that knowing facts is insufficient for avoiding falsehoods.

This matters for practitioners evaluating models. Selecting a model based on one truthfulness benchmark may lead to deployment failures on another dimension entirely. Yet no study has systematically quantified whether TruthfulQA, HaluEval, and FactScore measure the same underlying construct or tap distinct reliability dimensions.

The problem runs deeper than benchmark selection. At the surface level, researchers acknowledge that different benchmarks test different aspects of truthfulness. TruthfulQA uses multiple-choice questions targeting popular misconceptions; HaluEval employs binary detection of hallucinated details in QA, dialogue, and summarization; FactScore decomposes long-form generation into atomic facts verified against retrieval sources. These methodological differences are well-documented.

What remains unknown is whether these paradigmatic differences translate into genuinely independent capability dimensions. Do models that resist misconceptions also maintain generation coherence? Does factual precision in long-form generation correlate with either? Prior meta-evaluation work---BenchBench \cite{perlitz2024benchbench} and Ailem et al.\ \cite{ailem2024robustness}---established that benchmark agreement is not guaranteed and that methodological choices influence rankings. However, these studies examined general benchmark methodology, not truthfulness-specific correlation structure.

The gap is consequential. Without empirical correlation data, we cannot determine whether a single benchmark suffices for comprehensive truthfulness evaluation or whether multiple benchmarks capture complementary information. Model developers lack guidance on which dimensions require independent intervention.

We hypothesize that truthfulness in LLMs is a multi-dimensional construct comprising at least three distinct dimensions: misconception resistance (measured by TruthfulQA), generation coherence (measured by HaluEval), and factual precision (measured by FactScore). This hypothesis predicts moderate inter-benchmark correlations---higher than unrelated-benchmark baselines but substantially below unity---and a multi-factor structure in principal component analysis.

Our key insight is that different evaluation paradigms target different stages of the generation process. TruthfulQA tests whether models resist reproducing misconceptions encountered during training---a knowledge-selection problem. HaluEval tests whether models maintain consistency during text generation---a coherence problem. FactScore tests whether atomic claims in long-form outputs can be verified against external sources---a precision problem. A model can fail at any stage independently.

Building on this insight, we make the following contributions:

\begin{enumerate}
\item \textbf{First systematic cross-benchmark correlation analysis for truthfulness benchmarks.} We compute Spearman correlations between TruthfulQA, HaluEval, and FactScore across N=50 diverse LLMs spanning 7 architectures, 4 scales, and 4 training variants.

\item \textbf{Quantitative evidence for three-dimensional truthfulness.} We find inter-benchmark correlations of r=0.42--0.58, all exceeding the unrelated-benchmark baseline (r=0.10) but falling well below unity. PCA confirms 3 components are required to explain 80\% of variance; no single factor dominates (PC1 explains only 38.6\%).

\item \textbf{Divergent profile analysis.} We identify 4 models (8\%) exhibiting high MMLU but low TruthfulQA scores, demonstrating that general knowledge retrieval does not guarantee misconception resistance.

\item \textbf{Actionable evaluation guidance.} Our findings establish that single-benchmark evaluation is insufficient; comprehensive truthfulness assessment requires evaluating all three dimensions.
\end{enumerate}

The remainder of this paper is organized as follows. Section~\ref{sec:related} reviews related work on LLM evaluation and benchmark methodology. Section~\ref{sec:methodology} describes our correlation analysis methodology. Section~\ref{sec:experiments} details experimental setup across four hypotheses. Section~\ref{sec:results} presents results. Section~\ref{sec:discussion} discusses implications and limitations. Section~\ref{sec:conclusion} concludes.
